\Subhead{1.2.3}{Методы расчёта автоколебательных систем}
\textbf{Что ищем}: условие запуска, стационарные амплитуду и частоту, устойчивость и установление. Методы различаются тем, какие из этих задач они решают.
\begin{enumerate}
\item \textbf{Энергетический метод.} При почти гармонической форме считаем средний приток и потери. Их равенство даёт амплитуду, наклон разности --- её устойчивость. В первом приближении частота берётся близкой к линейной. Пример --- вывод \eqref{eq:vdpslow}.
\item \textbf{Гармонический баланс, квазилинейный расчёт.} Сохраняем основные гармоники нелинейного элемента. Комплексное усиление по петле $L(A,\omega)$ должно удовлетворять $L=1$: модуль задаёт амплитудный баланс, аргумент --- частоту. Порог находят из линейной модели малого поля; простое правило $|L|>1$ имеет смысл на частоте подходящей фазы и в выбранной модели. Само равенство $L=1$ устойчивости не доказывает.
\item \textbf{Метод медленно меняющихся амплитуд (МММА).} Вводим $A(t),\phi(t)$, усредняем по быстрому периоду; получаем укороченные дифференциальные уравнения. Они дают и стационарные решения, и возврат к ним, и поправку к частоте (см.~\Qref{1.4.2}).
\item \textbf{Качественный и численный анализ.} Для релаксационных и сильно несинусоидальных движений строят фазовый портрет, отображение Пуанкаре, интегрируют исходные уравнения. Один гармонический член уже не воспроизводит резкие скачки.
\end{enumerate}
Методы дополняют друг друга: баланс сначала объясняет, откуда берётся амплитуда, МММА показывает её эволюцию, фазовый анализ --- возможные режимы и их области притяжения.
